% Generated draft: requires mathematical and source review.
\documentclass[11pt]{article}
\usepackage[margin=25mm]{geometry}
\usepackage{fontspec}
\setmainfont{Latin Modern Roman}
\usepackage{amsmath,amssymb,mathtools,unicode-math}
\setmathfont{Latin Modern Math}
\usepackage{enumitem,needspace,xurl,hyperref}
\setlength{\parindent}{0pt}
\setlength{\parskip}{5pt}
\setlength{\emergencystretch}{3em}
\newcommand{\N}{\mathbb{N}}
\newcommand{\Z}{\mathbb{Z}}
\newcommand{\Q}{\mathbb{Q}}
\newcommand{\R}{\mathbb{R}}
\newcommand{\C}{\mathbb{C}}
\newcommand{\F}{\mathbb{F}}
\newcommand{\A}{\mathbb{A}}
\newcommand{\PP}{\mathbb{P}}
\newcommand{\E}{\mathbb{E}}
\newcommand{\eps}{\varepsilon}
\newcommand{\dd}{\,\mathrm{d}}
\begin{document}
\section*{UnsolvedMath ID: 30002382}
\textbf{Source problem code:} catalogue entry
\subsection{Definitions and mathematical statement}
An imaginary quadratic field is a degree-two extension \(K/\Q\) having no real embedding, considered up to \(\Q\)-isomorphism. Write \(D_K<0\) for its fundamental discriminant and \(h_K=|\operatorname{Cl}(\mathcal O_K)|\) for its ideal class number. Here prime class number means that \(h_K\) is a positive rational prime, so class number \(1\) is excluded. Let \(K^{\mathrm{ab}}\) be the maximal abelian extension of \(K\) inside a fixed algebraic closure, and set
\[
G_K^{\mathrm{ab}}=\operatorname{Gal}(K^{\mathrm{ab}}/K),
\]
with its Krull topology. Define
\[
\widehat{\Z}=\varprojlim_{m\geq 1}\Z/m\Z,
\qquad
A_0=\widehat{\Z}^{\times}\times\prod_{n\geq 1}\Z/n\Z,
\]
where the inverse limit is over divisibility, \(\widehat{\Z}^{\times}\) has its profinite topology, and the second factor has the product topology with each finite cyclic group discrete.

The supplied conjecture asks whether \(G_K^{\mathrm{ab}}\) is isomorphic to \(A_0\) for density one of the imaginary quadratic fields whose class number is prime. The natural profinite interpretation of isomorphism is a continuous group isomorphism with continuous inverse; the record does not explicitly specify this category.

For precision, a discriminant-ordered interpretation can be written as follows. For real \(X>0\), put
\[
\mathcal F(X)=\{K/\Q\text{ imaginary quadratic up to isomorphism}: |D_K|\leq X,\ h_K\text{ prime}\}.
\]
Under this ordering convention, the requested assertion is
\[
\lim_{X\to\infty}
\frac{\#\{K\in\mathcal F(X):G_K^{\mathrm{ab}}\cong_{\mathrm{top}}A_0\}}
{\#\mathcal F(X)}=1,
\]
with the quotient taken where its denominator is nonzero. This is relative density within the prime-class-number family. The source record does not specify an ordering or density convention, so this displayed limit is an explicit possible interpretation requiring source confirmation. The record also does not establish that the prime-class-number family is infinite; no such infinitude is asserted here.
\subsection{Short English statement}
The supplied conjecture predicts that almost all imaginary quadratic fields with prime ideal class number have the same absolute abelian Galois group, namely \(A_0\). A precise possible reading uses relative density ordered by absolute discriminant and isomorphisms of profinite groups. These conventions require confirmation from the source.
\subsection{Sources}
\begin{itemize}
\item Local UnsolvedMath registry, supplied record with permanent ID \(30002382\), title \emph{Abelian Galois Groups of Imaginary Quadratic Fields}, exact source problem code \texttt{catalogue entry}, in the Oberwolfach Reports Open Problems collection. The record gives proposed year \(2013\) and names no proposer.
\item Supplied report citation: \emph{Explicit Methods in Number Theory (2014)}, Oberwolfach Reports, pp.\ 2035--2082, DOI \texttt{10.4171/owr/2013/35}. Source URL: \url{https://doi.org/10.4171/owr/2013/35}. The supplied background describes Oberwolfach Reports as documenting mathematical workshops and their open-problem discussions. The citation's year \(2014\) and the report identifier \(35/2013\) are retained as supplied.
\item Supplied literature source: \emph{Imaginary quadratic fields with isomorphic abelian Galois groups}, \url{https://arxiv.org/abs/1209.6005}.
\item Supplied literature source: \emph{Oberwolfach Report 35/2013}, \url{https://doi.org/10.4171/owr/2013/35}.
\item Literature assessment dated exactly \texttt{2026-08-21}: the supplied curation labels the problem open and states that the universal-group description is proved for many imaginary quadratic fields, is supported numerically for prime class number, and that the density-one assertion has no proof. This is supplied dated curation derived from the OWR material, not an independently verified literature assessment; its status and cited claims require independent verification.
\item The registry labels the statement \texttt{corrected\_verified} and reports manual semantic, typographic, and LaTeX review of a meaning-preserving cleanup. These are supplied editorial assertions, not independent mathematical verification of this draft.
\end{itemize}
\end{document}
